\begin{frame}{Metodologia de validação}
  \small
  \begin{columns}[T,onlytextwidth]
    \begin{column}{0.31\textwidth}
      \begin{block}{1. Análise teórica}
        \begin{itemize}
          \item impedância $Z_C=1/(sC)$;
          \item aplicação da lei das correntes;
          \item resolução simbólica;
          \item projeto e diagramas de Bode.
        \end{itemize}
      \end{block}
    \end{column}
    \begin{column}{0.31\textwidth}
      \begin{block}{2. Simulação}
        \begin{itemize}
          \item três configurações no LTspice;
          \item fonte AC unitária;
          \item varredura logarítmica;
          \item buffer alimentado em $\pm10\,\mathrm{V}$.
        \end{itemize}
      \end{block}
    \end{column}
    \begin{column}{0.31\textwidth}
      \begin{block}{3. Ensaio em bancada}
        \begin{itemize}
          \item aferição dos componentes;
          \item montagem em protoboard;
          \item medição de $V_i$ e $V_o$;
          \item cálculo de $|H|=V_o/V_i$.
        \end{itemize}
      \end{block}
    \end{column}
  \end{columns}

  \begin{alertblock}{Condição de comparação}
    A tensão de entrada é medida em cada frequência.
    O corte corresponde a uma queda de $3\,\mathrm{dB}$ em relação ao ganho
    de baixa frequência de cada configuração.
  \end{alertblock}
\end{frame}
